%%%%%%%%%%%%%%%%%%%%%%%%%%%%%%%%%%%
% section 1 
\section{Introducing GL Controller}

\begin{frame}{1.0: GL Chamber: Introduction}{} 
GL Chamber will eventually replace P300-based (and SCP-based) chambers. The Platinous Series based on GL controller has been released. Work on Global chamber is underway. 
\begin{figure} [hbt!] 
\centering
\includegraphics[scale=0.26]{figures/glc-chamber-001a2.png}
\caption{ESPEC Platinous Series (EPL, EPU, EPX, EPX, J-series)}
\end{figure}
\end{frame}

\begin{frame}{1.1: What about Benchtop Chamber?}{} 
Benchtop chamber still uses Watlow F4T. Reason: Benchtop E-box is too small to fit GL controller hardware (I/O board, etc). For now, Benchtop chambers will be offered along with GL chambers. 
\begin{figure} [hbt!] 
\centering
\includegraphics[scale=0.32]{figures/benchtop-f4t.jpg}
\caption{E-box of Benchtop with ESPEC Web Controller}
\end{figure}
\end{frame}

\begin{frame}{1.2: Good News \& Bad News?}{}

\begin{minipage}[c]{0.50\textwidth}
{\bf{Good News?}} 
 We still offer our legacy ESPEC Web Controller (EWC) system for Benchtop chambers and U-Web kit for existing ESPEC (or Non-ESPEC with F4T) chambers that Customers still own. We can still apply our current knowledge of EWC to support customers.    
\end{minipage} 
\hfill
\begin{minipage}[c]{0.450\textwidth} 
        \includegraphics[width=\textwidth]{figures/benchtop-chambers-002.png} 
        \captionof{figure}{Benchtops with F4T}
\end{minipage}

\vspace{2em}
\begin{minipage}[c]{1.00\textwidth}
{\bf{Bad News?}} 
 We now have to learn about the GL chamber and its software to help potential customers. {\it{This short training will hopefully bring you up to date about the GL controller system.}} The ``good news'' part: GL controller in many ways is similar to our EWC. {\it{Details are explained in this presentation.}}  
\end{minipage}
\hfill
%\begin{minipage}[c]{0.450\textwidth} 
%        \includegraphics[width=\textwidth]{figures/up2board-io-002.PNG} 
%        \captionof{figure}{Ethernet ports on Up2board}
%\end{minipage}
\end{frame}

\begin{frame}{1.3: What about T-series Chambers?}{} 

\begin{minipage}[c]{0.550\textwidth}
{\bf{Good News?}} 
 T-series chamber configuration remains unchanged (still using Allen Bradley PLC). UI \& API of GL system will be ported to EWC to support T-series for improved performance.  
 
 \vspace{1em} We still offer our legacy ESPEC Web Controller (EWC) system for T-series chambers and U-Web kit for existing T-series chambers customers still own. We can still apply our current knowledge of EWC to support customers.    
\end{minipage} 
\hfill
\begin{minipage}[c]{0.4250\textwidth} 
        \includegraphics[width=\textwidth]{figures/t-series-001.png} 
        \captionof{figure}{T-series with EWC}
\end{minipage}

%\vspace{2em}
%\begin{minipage}[c]{1.00\textwidth}
%{\bf{Bad News?}} 
% We now have to learn about the GL chamber and its software to help potential customers. {\it{This short training will hopefully bring you up to date about the GL controller system.}} 
%\end{minipage}
%\hfill
%\begin{minipage}[c]{0.450\textwidth} 
%        \includegraphics[width=\textwidth]{figures/up2board-io-002.PNG} 
%        \captionof{figure}{Ethernet ports on Up2board}
%\end{minipage}
\end{frame}

%%%%%%%%%%%%%%%%%%%%%%%%%%%%%%%%%%%
% section 2
\section{What is GL Controller?}

\begin{frame}{2.0: GL Controller: What is GL Chamber?}{}
\begin{minipage}[c]{0.57\textwidth}
GL controller is part hardware and part software that makes up the ``GL chamber." 

\vspace{0.1in} 
The hardware part makes up the Input/Output (I/O) board (see figure).

\vspace{0.1in}
The software part is written in CODESYS; its details do not concern us! This software part works with Up2Board computer to allow GL to operate and control its I/O ports. This Up2Board also supports our EWC, as depicted in the figure below. 
\begin{figure} [hbt!] 
      \centering
      \includegraphics[scale=.25]{figures/up2board-001a.PNG}  
      %\caption{EWC V3 log-in screen on a Web browser}
\end{figure}
How does GL controller system work? 
\end{minipage} 
\hfill
\begin{minipage}[c]{0.410\textwidth} 
        \includegraphics[width=\textwidth]{figures/platinous-jseries-001.png} 
        \captionof{figure}{GL software/hardware in Platinous J Series} 
\end{minipage}
\end{frame} 

\begin{frame}{2.1: GL System vs. Legacy EWC}{}
\begin{minipage}[c]{0.520\textwidth}
Figure on the right depicts how a Benchtop chamber and EWC work together to provide control and operation of the chamber. Chamber interface and communication between F4T and EWC are via Ethernet or serial (RS232) connection. Depicted in the figure, two Ethernet cables run from EWC, one to F4T, the other to external Ethernet port to join the main network for remote access of the EWC user interface. 

\vspace{1em}
In contrast, GL system and EWC are integrated, to be described shortly. 

\end{minipage} 
\hfill
\begin{minipage}[c]{0.4250\textwidth} 
    \includegraphics[width=\textwidth]{figures/benchtop-f4t.jpg}
    \captionof{figure}{EWC and Benchtop F4T} 
\end{minipage}
\end{frame} 

\begin{frame}{2.2: Details of EWC and P300 Interface (Review)}{}
\begin{minipage}[c]{0.580\textwidth}
Figure on the right depicts communication interface between EWC and P300. We already know, EWC connects to a P300 controller via a serial interface to provide control and operation of the chamber.

\vspace{1em}
The same communication protocol is also used on SCP-220 (or ES102) or Watlow F4. 

\vspace{1em}
{\bf{Summary:}} Legacy EWC functions as a separate device to connect to a controller (such as P300, SCP220, F4/F4T, etc) to provide control and operation of the chamber. 
\end{minipage} 
\hfill
\begin{minipage}[c]{0.380\textwidth} 
  \includegraphics[width=\textwidth]{figures/ewc-n-p300-01a.PNG}
  \captionof{figure}{P300 Chamber} 
\end{minipage}
\end{frame} 

\begin{frame}{2.3: Details of EWC and T-series (Review)}{}
T-series and EWC are communicated through an internal network via ETH1, while external network is established between EWC and customer main netowrk via ETH0 to allow access to control and operate the chamber remotely. 

  \begin{figure} [hbt!] 
      \centering
      \includegraphics[width=0.85\textwidth]{figures/tseries-n-ewc.PNG}
      \caption{Chamber interface configuration between T-series and EWC} 
  \end{figure}

\end{frame} 

\begin{frame}{2.4: Detailed Configuration of EWC \& T-series}{}
EWC box is mounted on top of T-series chamber with Ethernet cables running through the E-box, with ETH-1 connecting to the PLC, ETH-0 to the chamber Ethernet port that connects to the outside world. 

  \begin{figure} [hbt!] 
     \centering
    \includegraphics[width=0.70\textwidth]{figures/u-web-kit-001.PNG}
     \caption{Chamber interface configuration between T-series and EWC} 
  \end{figure}

\end{frame} 

\begin{frame}{2.5: Details of GL \& Controller Interface}{}
\begin{minipage}[c]{0.530\textwidth}

Contrary to the legacy EWC system, which offers chamber interface and configuration as an external connection (i.e. because EWC is implemented to support a variety of controllers), the GL controller has {\bf{no external chamber interface}}. 

\vspace{1em} GL controller {\it{is}} the CODESYS that resides on the Up2Board. This Up2Board and CODESYS are the GL controller. The GL controller uses Up2Board to link to its I/O board, depicted in figure on the right. 
\end{minipage} 
\hfill
\begin{minipage}[c]{0.420\textwidth} 
    \includegraphics[width=\textwidth]{figures/glc-chamber-003.png}
    \captionof{figure}{GL Hardware} 
\end{minipage}
\end{frame} 

\begin{frame}{2.6: Detailed Interface of GL Controller}{}
\begin{minipage}[c]{0.530\textwidth}

Interface \& I/O ports of Up2Board used by GL controller.

\begin{figure} [hbt!] 
\centering
  \includegraphics[scale=0.250]{figures/gl-up2board-001a.png}
  \caption{Interfaces used by GL Ctrl} 
\end{figure}
    
\begin{itemize} 
  \item [(1)] Supplies power to HMI 
  \item [(2)] Connects to HMI operation 
  \item [(3)] links CODESYS license; fourth USB for external storage
  \item [(4)] Connects to HDMI display 
  \item [(5)] Connects to main network  
\end{itemize} 
\end{minipage} 
\hfill
\begin{minipage}[c]{0.420\textwidth} 
        \includegraphics[width=\textwidth]{figures/glc-chamber-004.png}
        \captionof{figure}{GL Chamber} 
\end{minipage}
\end{frame} 

\begin{frame}{2.7: Chamber User Interface (HMI)}{}
Complete control and operation of the chamber are accomplished via the dedicated HMI on the chamber. {\it{GL user interface explorered later.}}

    \begin{figure} [hbt!] 
      \centering
      \includegraphics[scale=0.250]{figures/gl-hmi-003a.png}
      \caption{GL Touchscreen Display, dedicated HMI} 
    \end{figure}
\end{frame}   

\begin{frame}{2.8: Details of External Ports of GL Chamber}{}

External ports of the GL controller (Platinous J Series) has three ports: (1) Product Temp. Protector, (2) Chamber light switch, (3) External Storage (USB Thumb drive for downloading data, uploading update packages, etc). 

    \begin{figure} [hbt!] 
      \centering
      \includegraphics[scale=0.40]{figures/gl-external-ports-001a.png}
      \caption{External ports on GL Platinous J series chamber} 
    \end{figure}
\end{frame}   

%%%%%%%%%%%%%%%%%%%%%%%%%%%%%%%%%%%
% section 3
\section{GL Software Update and Distribution} 
\begin{frame}{3.0: GL Software Update}{}
Similar to our current (legacy) EWC, GL system uses chamber S/N as its registration number to receive firmware update through our cloud service provider, called {\it{MenderIO}}.
    \begin{figure} [hbt!] 
      \centering
%      \includegraphics[scale=.30]{figures/menderio-cloud-001.PNG}  
      \includegraphics[scale=0.300]{figures/mender-cloud-service-002.PNG}
      \caption{GL System subscribes to MenderIO cloud service}
    \end{figure}
\end{frame} 

\begin{frame}{3.1: Firmware Update: Online or Offline}{} 
Customer has complete control over firmware update method: Offline or Online through MenderIO cloud service. 
\begin{enumerate}
\item Using a stepladder technique, update package is pushed onto the next available root partition (see diagram). 

    \begin{figure} [hbt!] 
      \centering
%      \includegraphics[scale=.30]{figures/menderio-cloud-001.PNG}  
      \includegraphics[scale=.60]{figures/stepladder-001.PNG}
      \caption{Stepladder technique of GL update, (same as EWC)}
    \end{figure}

\item If the update process is successful, the system will boot and run on this new firmware.   
\item If the update process failed, the system rolls back to the previous known operation state. 
\end{enumerate} 
\end{frame} 

\begin{frame}{3.2: Firmware Update Subscription}{} 
Firmware update is provided through a yearly update subscription. Customers purchase the subscription 
to receive the updates. 
\begin{itemize}
   \item {\bf{Online Update}}: \\
       For GL chamber that connects to the Internet, MenderIO cloud updating service handles the update process. 
       Automatic update will take place when GL chamber is in Standby mode.   
   \item {\bf{Offline Update}}:   \\
       For GL chamber that does not connect to the Internet, the owner is provided with a link to access their account
       to download an update package to perform the update procedure on their system. Refer to User's
       Manual for detail on Firmware Update procedure. 
\end{itemize} 
\end{frame} 

%%%%%%%%%%%%%%%%%%%%%%%%%%%%%%%%%%%
% section 4
\section{User Interface: Display Options, Features \& Operation}

\begin{frame}{4.0: GL Operation Panel}{}

The operating panel is the 10-inch touchscreen HMI that provides the User Interface (UI) for control and operation of the chamber. Depicted in the figure below, the UI consists of three major components: (1) Menu bar, (2) Status bar, and (3) Standard Display area. The touchscreen is operated by gently pressing the screen elements provided by these components.
\begin{figure} [hbt!] 
      \centering
      \includegraphics[scale=.34]{figures/gl-ui-001.PNG}  
      \caption{User Interface of GL Controller}
\end{figure}
\end{frame} 

\begin{frame}{4.1: GL Main Menu \& Features}{} 
The menu bar (left side) provides access to the GL controller system via five separate menus. 
\begin{enumerate}
   \item {\bf{Home}}: GL controller Overview page. Displays menu buttons, custom buttons and chamber operation status.
   \item {\bf{Operation Mode}}: Controls Constant (1, 2, 3) modes and Program profiles.    
   \item {\bf{Monitor}}: Monitor GL chamber operating status; five subpages are available to manage the monitoring status
   \item {\bf{Set Mode}}: Configure Constant Mode settings; create or edit program profiles.
   \item {\bf{Setup}}: System-wide configuration and settings
\end{enumerate} 

Hands-on will provide a better understanding of these menus; many of which are similar to our current EWC menu and user interface. 

\end{frame}

\begin{frame}{4.2: Three Ways to Access GL Menu}{}

GL Controller system offers three ways to access its system menu. 

\begin{figure} [hbt!] 
      \centering
      \includegraphics[scale=.3250]{figures/three-access-way-001.PNG}  
      \caption{User Interface of GL Controller}
\end{figure}

\begin{itemize} 
    \item[(1)] Menu bar 
    \item[(2)] Menu Button (via the Home page)
    \item[(3)] Hamburger Button (in Status bar) 
\end{itemize} 
\end{frame} 


%%%%%%%%%%%%%%%%%%%%%%%%%%%%%%%%%%%
% section 5
\section{GL Chamber \& Network Configuration} 
\begin{frame}{5.0: Nework Configuration: DHCP or Static}{}
\begin{itemize} 
\item By default, the GL controller operates as a standalone system. Complete control and operation of the instrumentation are achieved via the dedicated HMI on the chamber. 

\item When the GL controller is connected to a network, you can use a Web browser to access the GL controller user interface (UI) to control and operate the chamber. The GL controller utilizes a standard Web browser to provide (i.e., host) its UI for control and operation of the instrumentation. Authorized users can access and control the GL controller/chamber remotely. 

\item Network connection and protocol are standard (DHCP or Static) based on our legacy network protocol on ESPEC Web Controller. However, remote access to the GL controller on the network must be enabled on the HMI via the Set Protection menu under {\bf{Setup}} (details outlined below). 
\end{itemize} 
\end{frame}

\begin{frame}{5.1: How to Connect GL Controller to the Main Network} {} 
\begin{minipage}[c]{0.550\textwidth}
The GL controller system connects to a local network with the Ethernet cable running from the network port in the back of the chamber marked "Web Controller" (see figure on the right) to the main network Ethernet port. 

\vspace{1em} 
By default, the GL controller system uses a dynamic host control protocol (DHCP) to join any network.  
\end{minipage} 
\hfill
\begin{minipage}[c]{0.40\textwidth} 
        \includegraphics[width=\textwidth]{figures/gl-network-connect-001.PNG}
        \captionof{figure}{GL network connection}
\end{minipage}

\end{frame}

\begin{frame}{5.2: DHCP Network Configuration} {} 
As part of a network device, the GL controller system can still operate as a standalone system (with its remote access feature disabled by default). A unique IP address assigned by the DHCP server will be used by the GL controller system. This IP address can be found under the {\bf{Setup}} menu in the {\bf{Set Communication}} page (under {\bf{Configuration}}), see next slide. 
       \begin{figure} [hbt!] 
         \centering
         \includegraphics[scale=0.55]{figures/gl-dhcp-ip.PNG}  
         \caption{DHCP network diagram}
       \end{figure}
\end{frame}

\begin{frame}{5.3: Accessing GL Controller on a DHCP Network} {} 
Remote access to the GL controller is via its hostname or its IP address:
\begin{itemize}
    \item {\url{http://hostname.local/}} \\ Example (see arrow 1): {\url{http://GL-EPL-3-P300.local/}} 
    \item {\url{http://IP-address/}} \\ Example (see arrow 2): {\url{http://10.30.100.190/}}
\end{itemize} 
       \begin{figure} [hbt!] 
         \centering
         \includegraphics[scale=0.20]{figures/glc-dhcp-ip-addrs.PNG}  
         \caption{Accessing GL controller via a network connection}
       \end{figure}
\end{frame}

%\begin{frame}{5.4: DHCP by Default} {} 
%By default, the GL controller system uses a dynamic host control protocol (DHCP) to join any network.  
%       \begin{figure} [hbt!] 
%         \centering
%         \includegraphics[scale=0.20]{figures/glc-dhcp-ip-addrs.PNG}  
%         \caption{DHCP network diagram}
%       \end{figure}
%\end{frame}

\begin{frame}{5.4: Static Network Configuration} {} 
A static IP address can be used for the GL controller system to join a network. This setup requires: (1) IP address, (2) Net Mask, (3) Gateway, (4) DNS 1 and DNS 2. These parameters must be configured manually.  

\vspace{1em}
%      \begin{figure} [hbt!] 
%        \centering
%        \includegraphics[scale=0.70]{figures/gl-static-general.PNG}  
%        \caption{General static network diagram}
%      \end{figure}

\begin{minipage}[c]{0.35\textwidth}
  Example: 
    \begin{itemize} 
        \item {\it{IP}}: 10.30.200.247 
        \item {\it{Mask}}: 255.255.0.0 
        \item {\it{Gateway}}: 10.30.0.1
        \item {\it{DNS1}}: 10.30.30.31 
        \item {\it{DNS2}}: 8.8.8.8
    \end{itemize} 
\end{minipage} 
\hfill
\begin{minipage}[c]{0.63\textwidth} 
        \includegraphics[width=\textwidth]{figures/gl-static-general.PNG}
        \captionof{figure}{GL Static Network} 
\end{minipage}

\end{frame}

\begin{frame}{5.5: Static Network Configuration} {} 
Here is an example of a static network configuration found on the Set LAN page of the GL controller on the HMI. Key points for custom static configuration are items (1) and (2). DHCP box must be unchecked and complete network configuration parameters must be recorded in the five fields under item (2), as shown. 
      \begin{figure} [hbt!] 
        \centering
        \includegraphics[scale=0.250]{figures/glc-static-network-001.PNG}  
        \caption{General static network configuration}
      \end{figure}
\end{frame}

\begin{frame}{5.6: Static/Fallback Network Configuration} {} 
It is easy to use GL controller on a small (or private) network without DHCP service. The GL controller can connect directly to a client PC (see diagram) or be part of a small network. It will use a preconfigured IP setting based on a Class C network protocol. This IP address (192.168.0.83) is called {\it{fallback IP}}.  

\begin{minipage}[c]{0.35\textwidth}
  Example: 
    \begin{itemize} 
        \item {\it{IP}}: 192.168.0.83 
        \item {\it{Mask}}: 255.255.255.0 
        \item {\it{Gateway}}: 192.168.0.1
        \item {\it{DNS1}}: 192.168.0.1 
        \item {\it{DNS2}}: 8.8.8.8
    \end{itemize} 
\end{minipage} 
\hfill
\begin{minipage}[c]{0.63\textwidth} 
        \includegraphics[width=\textwidth]{figures/gl-fallback-ip.PNG}
        \captionof{figure}{GL Private network diagram} 
\end{minipage}

      %\begin{figure} [hbt!] 
      %  \centering
      %  \includegraphics[scale=0.60]{figures/gl-fallback-ip.PNG}  
      %  \caption{Private network diagram}
      %\end{figure}
\end{frame}

\begin{frame}{5.7: Fallback IP Address of GL Controller} {} 
When DHCP service is not detected, GL controller will use its preconfigured IP address (fallback IP). The DHCP box in the figure below ensures that if a DHCP service is present, GL controller will request an IP address to join the network.  Only when that attempt fails does it ``fall back'' to use its preconfigured IP address.  
      \begin{figure} [hbt!] 
        \centering
        \includegraphics[scale=0.280]{figures/gl-fallback-cfg-001a.PNG}  
        \caption{Private network diagram}
      \end{figure}
\end{frame}

\begin{frame}{5.8: GL Chamber: Standalone}{}

GL Chamber can operate by itself without connecting to a network or an Internet. It can still operate as a standalone system even when it is part of a network.  
\begin{figure} [hbt!] 
      \centering
      \includegraphics[scale=.30]{figures/glc-chamber-001a2.png}
      \caption{GL a standalone system}
\end{figure}
\end{frame} 

\begin{frame}{5.9: GL Chamber: Network Communication}{}

As part of a network, GL chamber can be accessed remotely using a Web browser via IP address or hostname of the GL chamber. However, remote access must be enabled for this to happen. It must be enabled on the HMI via the Setup submenu ``Set Protection''. This option is linked to remote communication in the next slide. 

\begin{figure} [hbt!] 
      \centering
      \includegraphics[scale=0.40]{figures/gl-network-001.png}  
      \caption{GL chamber as part of a network device}
\end{figure}
\end{frame} 

%%%%%%%%%%%%%%%%%%%%%%%%%%%%%%%%%%%
% section 6
\section{GL Remote Control \& ChamberConnectLibrary}

\begin{frame}{6.0: GL Remote Communication}{} 
Communication protocol of the GL controller system is a general-purpose communication tool that allows communication between the GL chamber and remote device (PC). Communication protocol is currently available for TCP/IP and RS232 interface, with TCP/IP set by default. This option allows customers to control the GL chamber remotely (or locally) using GL native text commands. 

\vspace{1em} 

\begin{itemize} 
  \item {\bf{Locally}}: This feature is available under the Set Communication submenu: {\keys{Setup}}$\longrightarrow${\keys{Configuration}}$\longrightarrow${\keys{Set Communication}}
  \item {\bf{Remotely}}: Remote communication can be enabled/disabled. This option is available under the Set Protection submenu: {\keys{Setup}}$\longrightarrow${\keys{Set Protection}} or Set Protection button on HMI.   
\end{itemize}
The next slide depicts enable/disable options for remote communication and Web access. Hands-on demo will be provided later.   
\end{frame}

\begin{frame}{6.1: Setting Remote Access on GL}{} 
Remote access (communication and Web) must be enabled on the HMI.  
\begin{figure} [hbt!] 
      \centering
      \includegraphics[scale=0.22]{figures/set-protection-remote-comm-001.png}  
      \caption{Disable/enable GL remote access for COMM and Web}
\end{figure}
\end{frame} 

\begin{frame}{6.2: GL \& ChamberConnectLibrary}{} 
\begin{itemize} 
    \item Chamber Connect Library using Python 3 will be available for customers who prefer to customize their programs to control their GL chamber remotely via Python 3 programming. \\

    \vspace{2em} 

    \item The Library will be integrated with the current Chamber Connect Library for P300, SCP220, ES102 and Watlow F4 and F4T. \\ {\bf{Note:}} T-series PLC is unsupported. \\

    \vspace{2em} 

    \item Communication protocol and interface for GL include TCP/IP and Serial. GP-IB communication protocol is currently not supported. 
\end{itemize} 
\end{frame}

\begin{frame}{6.3: GL \& ChamberConnectLibrary Support Page}{} 
Customer support page and ChamberConnectLibrary link to GitHub. 
\begin{figure} [hbt!] 
      \centering
      \includegraphics[scale=0.28]{figures/gl-chamberconnectlib-py3-02.png}  
      \caption{Support page for Chamber Conenct Library}
\end{figure}

\end{frame}
%%%%%%%%%%%%%%%%%%%%%%%%%%%%%%%%%%%
% section 7
\section{Troubleshooting GL Chamber} 

\begin{frame}{7.0: Troubleshooting Issues on GL}{}

Integration of hardware \& software of the GL controller has a number of advantages. Issue(s) customer may have is virtually not about chamber misconfiguration or chamber connection interface any more. How many times have we helped customers with their chamber interface connection and misconfiguration that they themselves did.  \newline 

\begin{enumerate} 
\item Server Error /Chamber error: Stopped Working
\item Unable to access GL chamber remotely
\item Unable to issue set commands directly to GL chamber
\item Network connection issues 
\end{enumerate} 
\end{frame}  

%%%%%%%%%%%%%%%%%%%%%%%%%%%%%%%%%%%
% section 8
\section{Universal Web Controller Kit for GL Chamber?} 

\begin{frame}{8.0: U-Web Kit Support: COMM Interface}{} 

Offer of the Universal Web Controller kit (U-Web kit) continues with its standard configuration options to provide connection, operation and control on existing ESPEC or Non-ESPEC (with F4T) chambers. We still do offer this option. Default communication with F4T is via TCP/IP; with other chambers/controllers, it will use the serial COMM via the Serial-to-USB interface, as shown below.  

    \begin{figure} [hbt!] 
      \centering
      \includegraphics[scale=.220]{figures/webkit0.PNG}  \\
      \caption{U Web kit hardware in box and bare}
    \end{figure}
\end{frame}

\begin{frame}{8.1: U-Web Kit for Legacy Chambers}{} 
EWC can be sold as a U-Web Kit to connect and control existing ESPEC chambers or non-ESPEC chambers (mentioned above). 

    \begin{figure} [hbt!] 
      \centering
      \includegraphics[scale=.260]{figures/u-web-kit-up2board-001.PNG}  
      \caption{U-Web Kit with USB-to-Serial interface cable}
    \end{figure}
{\bf{Advantages:}} 
\begin{itemize}
\item U-Web kit is an excellent option for adding EWC to the existing chamber for control and operation. 
\item Setup is generally made simple with its quick guide and preconfiguration. 
\end{itemize} 
\end{frame}

\begin{frame}{8.2: U-Web Kit for GL Chambers}{} 
\begin{itemize} 
  \item {\bf{Question:}} Is there a U-Web kit for GL chambers? 
  \item {\bf{Answer:}} There is no U-Web kit for GL chambers. \\ GL {\it{is}} the Web Controller. 
\end{itemize} 
\end{frame}

%%%%%%%%%%%%%%%%%%%%%%%%%%%%%%%%%%%
% section 9
\section{Hands-On Demo (HMI) \& Questions/Answers} 

\begin{frame}{9.0: GL User Interface}{}
  
\begin{itemize}
    \item Quick overview of GL Operation page

    \begin{figure} [hbt!] 
      \centering
      \includegraphics[scale=.150]{figures/gl-chamber-op1.PNG}  
      \caption{Program/Preview menu of GL controller}
    \end{figure}
    \item Quick overview of GL Monitor page 
    \begin{figure} [hbt!] 
      \centering
      \includegraphics[scale=.150]{figures/gl-chamber-op2.PNG}  
      \caption{Program execution and monitoring on GL chamber}
    \end{figure}
\end{itemize}
\end{frame}

\begin{frame}{9.1: How to Log in via HMI}{}
  Users unaware of the special access menu may have trouble logging in. 
\begin{itemize}
    \item Login status: User logged out
    \begin{figure} [hbt!] 
      \centering
      \includegraphics[scale=.150]{figures/user-logout-001-40-54.PNG}  
      \caption{Home page of GL; not user logged in}
    \end{figure}
\end{itemize}
\end{frame}

\begin{frame}{9.2: How to Enter Password to Log in via HMI}{}
  Users unaware of the special access menu may have trouble logging in. 
\begin{itemize}
    \item Click on user icon (with diagonal bar) to log in, but cannot enter password
    \begin{figure} [hbt!] 
      \centering
      \includegraphics[scale=.150]{figures/auto-login-opt-41-08.PNG}  
      \caption{Require keybaord enable to enter password; how?}
    \end{figure}
    \item Answer: Review our GL Oepration Manual (illustration, Section 2.6.3)
\end{itemize}
\end{frame}
\begin{frame}{9.3: Q/A? \& Demo Units}{}
\begin{itemize}
   \item Questions \& Answers
   \item GL chamber demo: \\
\end{itemize}
   \vspace{1em} 
   GL R \& D demo unit on our LAN: \\
   {\url{http://10.30.200.247/}} \\
   {\it{username:}} cservice \\ {\it{password:}} cservice \\
   \vspace{1em}
   Demo unit for potential customers on a public network: \\
   Pulbic Netwrok: {\url{http://50.217.63.109/}} \\
   {\it{username:}} demo \\ {\it{password:}} demo \\
   \vspace{1em}  
   Private IP (on our LAN): {\url{http://192.168.250.4/}} \\
   {\it{username:}} demo \\ {\it{password:}} demo \\
   \vspace{1em}  
\end{frame}
